\documentclass{article}
\usepackage[UTF8]{ctex}
\usepackage{amsmath,amssymb,amsfonts}
\usepackage{enumitem}
\usepackage{geometry}
\geometry{a4paper, margin=1in}
% === 错题本排版优化（不参与题目标识） ===
\setlength{\parindent}{0pt}
\setlength{\parskip}{0.35em plus 0.10em minus 0.08em}
\setlength{\emergencystretch}{3em}
\setlist{itemsep=0.25em, topsep=0.40em, parsep=0pt}
\allowdisplaybreaks[2]
% === 错题本排版优化结束 ===
\begin{document}
\section*{郑州大学2025年数学分析真题}
\section*{128．郑州大学2025年数学分析试卷}

\subsection{题目1}
1．（10 分）求极限 \(\lim _{x \rightarrow 0}\left[\frac{1+\int_{0}^{x} \sin t \mathrm{~d} t}{\ln \left(1+x^{2}\right)}-\frac{1}{\sin ^{2} x}\right]\) ．

\subsection{题目2}
2．（10 分）求极限 \(\lim _{n \rightarrow \infty} \frac{\left(1+2^{a}+3^{a}+\cdots+n^{a}\right)^{b+1}}{\left[1+3^{b}+5^{b}+\cdots+(2 n+1)^{b}\right]^{a+1}}\) ，其中 \(a>-1, b>-1\) ．

\subsection{题目3}
3．（10 分）求不定积分 \(\int \frac{x^{2}}{(1-x)^{2024}} \mathrm{~d} x\) ．

\subsection{题目4}
4．（10 分）判断级数 \(\sum_{n=1}^{\infty}(-1)^{n-1}\left[e-\left(1+\frac{1}{n}\right)^{n}\right]\) 的绝对收敛性和条件收敛性．

\subsection{题目5}
5．（10 分）设 \(f(x)=\arctan 2 x\) ，求 \(f^{(n)}(0)\) ．

\subsection{题目6}
6．（15 分）设 \(D=[0,1] \times[0,1]\) ．

\begin{itemize}
  \item[（1）] 计算 \(A=\iint_{D}\left|x y-\frac{1}{4}\right| \mathrm{d} x \mathrm{~d} y\) ．
  \item[（2）] 设 \(z=f(x, y)\) 在 \(D\) 上连续，且满足 \(\iint_{D} f(x, y) \mathrm{d} x \mathrm{~d} y=0, \iint_{D} x y f(x, y) \mathrm{d} x \mathrm{~d} y=1\) ，证明：存在\(\left(x_{0}, y_{0}\right) \in D\) ，使得 \(\operatorname{Af}\left(x_{0}, y_{0}\right) \geq 1\) ．
\end{itemize}

\subsection{题目7}
7．（15 分）已知二元函数

\[
f(x, y)= \begin{cases}\frac{x^{3}}{x^{2}+y^{2}}, & (x, y) \neq(0,0) \\ 0, & (x, y)=(0,0)\end{cases}
\]

\begin{itemize}
  \item[（1）] 求 \(f_{x}(x, y), f_{y}(x, y)\) 。
  \item[（2）] 设方向 1 与 \(x\) 轴正方向的夹角为 \(\alpha\) ，求 \(\left.\frac{\partial f(x, y)}{\partial 1}\right|_{(0,0)}\) ．
  \item[（3）] 证明 \(f(x, y)\) 在点 \((0,0)\) 处连续但不可微．
\end{itemize}

\subsection{题目8}
8．（15 分）设 \(p>1, q \in \mathbb{R}\) ，证明：函数项级数 \(\sum_{n=1}^{\infty} \frac{1}{n^{p}+n^{q} x^{2}}\) 在 \((-\infty,+\infty)\) 上一致收敛，且当 \(q<3 p-2\)时可逐项微分。

\subsection{题目9}
9．（20 分）证明下列各题。

\begin{itemize}
  \item[（1）] （5 分）设 \(a\) 为大于 0 的常数，当 \(x>0\) 时，\(f(x)=\frac{1}{\sqrt{x}}-\frac{1}{\sqrt{x+a}}\) 单调递减。
  \item[（2）] （7 分） \(\int_{0}^{\sqrt{2 \pi}} \sin x^{2} \mathrm{~d} x>\frac{2-\sqrt{2}}{2 \sqrt{\pi}}\) ．
  \item[（3）] （8 分） \(\int_{0}^{+\infty} \sin x^{2} \mathrm{~d} x\) 条件收敛．
\end{itemize}

\subsection{题目10}
10．（20 分）设 \(f(x)\) 在 \([a, b]\) 上不恒为常数，在 \([a, b]\) 上二阶可导，且满足 \(f(a)=f(b), \lim _{x \rightarrow a^{+}} \frac{f(x)-f(a)}{(x-a)^{2}}=\) 2．证明：

\begin{itemize}
  \item[（1）] \(f_{+}^{\prime}(a)=0, f_{+}^{\prime \prime}(a)=4\) ．
  \item[（2）] 存在 \(\xi \in(a, b)\) ，使得 \(f^{\prime \prime}(\xi)=0\) ．
  \item[（3）] 存在 \(\eta \in(a, b)\) ，使得 \(f^{\prime}(\eta)=\frac{f(\eta)-f(a)}{\eta-a}\) ．
\end{itemize}

\subsection{题目11}
11．（15 分）设连续函数列 \(\left\{f_{n}(x)\right\}\) 在 \([a, b]\) 上一致收敛于 \(f(x)\) ，而 \(g(u)\) 在 \((-\infty,+\infty)\) 上连续，证明：\(\left\{g\left(f_{n}(x)\right)\right\}\) 在 \([a, b]\) 上一致收敛于 \(g(f(x))\) ．
\end{document}
